
Este capítulo retoma o percurso realizado na revisão para o TCC~II, confronta o que foi construído com os objetivos estabelecidos na introdução, explicita as limitações do que se apresenta e delimita os encaminhamentos previstos para a etapa seguinte.

\section{Retomada dos Objetivos}

O objetivo geral consistia em projetar e desenvolver um interpretador dinâmico, integrado a uma IDE web, que permitisse ao usuário personalizar a sintaxe e o vocabulário da linguagem. Ao término desta etapa, o núcleo de tradução encontra-se operante e integrado à plataforma, e os objetivos específicos apresentam a seguinte situação.

O primeiro objetivo, especificar a gramática da linguagem e os elementos passíveis de redefinição em tempo de execução, foi cumprido e está registrado na Seção~\ref{sec:especificacao-linguagem} e na Tabela~\ref{tab:elementos-configuraveis}, que relaciona cada elemento configurável, sua forma canônica e as restrições verificadas antes do início da análise.

O segundo e o terceiro, projetar e implementar os módulos de análise léxica, sintática e semântica com suporte à configuração de vocabulário e à seleção das variantes previstas na implementação, foram cumpridos. As fases encadeiam-se do analisador léxico à execução pelo interpretador, conforme detalhado na Seção~\ref{sec:metodologia-interpretador}, e a configuração da linguagem é recebida e validada a cada nova compilação.

O quarto, integrar o núcleo à interface web definindo os contratos de comunicação entre os módulos, foi cumprido por meio da organização em \textit{monorepo} descrita na Seção~\ref{sec:modelagem-arquitetura}, que permite à IDE consumir o pacote \texttt{compiler} como dependência local.

O quinto, disponibilizar o núcleo em ambiente web e verificar seu comportamento por meio de uma suíte de testes automatizados, foi cumprido no que diz respeito à disponibilização e à cobertura das construções da linguagem, conforme a Seção~\ref{sec:metodologia-testes}.

A revisão de 16 de setembro de 2026 acrescenta ao relato a igualdade estrita, a expressão condicional ternária, a validação de tipos durante a execução e o suporte à cedilha no léxico. A suíte do núcleo aprovou 224 casos em 21 arquivos, incluindo os comportamentos introduzidos. Esse resultado não implica compatibilidade integral com as linguagens que inspiram as predefinições da IDE.

A revisão de 19 de setembro de 2026 acrescenta a descrição dos contratos de resposta e dos controles do serviço Python. A execução local aprovou os 188 testes do back-end em 17 arquivos e confirmou os 224 casos do compilador. Os exemplos de minimização incluem resumos de linguagens, respostas de usuário sem senha e histórico de submissões restrito ao solicitante, com omissão do código enviado \cite{projeto2026backend,silva2026verificacaobackend}.

O sexto objetivo, definir o protocolo de avaliação da solução, foi atendido pela Seção~\ref{sec:protocolo-avaliacao}, que descreve quem responde, os blocos de itens, as escalas usadas e o cálculo dos escores, além dos requisitos já apresentados nas Tabelas~\ref{tab:requisitos-funcionais} e \ref{tab:requisitos-nao-funcionais}. Para aplicar o instrumento ainda falta definir a amostra e obter o encaminhamento de ética da instituição, e o Capítulo~\ref{cap:resultados} registra essa pendência. Vale separar o que o protocolo permite concluir: como é dirigido a um painel de especialistas, ele coleta a opinião de profissionais sobre a plausibilidade da contribuição pedagógica. Provar essa contribuição exigiria um estudo com estudantes.

\section{Limitações}

Cabe explicitar os limites do que aqui se apresenta, de modo que os resultados não sejam lidos além do que sustentam.

A limitação mais relevante diz respeito à ausência de verificação empírica. A pergunta que orienta a pesquisa indaga em que medida a personalização da sintaxe contribui para reduzir a barreira do primeiro contato com a programação; o que este trabalho entrega é o artefato que torna essa investigação possível, e não a resposta a ela. Nenhuma afirmação sobre ganho de aprendizagem pode ser sustentada com os dados atualmente disponíveis.

No plano técnico, a linguagem suportada é um subconjunto didático de linguagens imperativas, com declaração explícita de tipo ou modo dinâmico. Ela cobre as construções típicas das disciplinas introdutórias, mas não contempla orientação a objetos, tratamento estruturado de exceções nem modularização em múltiplos arquivos de código-fonte.

A personalização abrange o vocabulário, operadores em forma de palavra, literais booleanos e delimitadores, além da seleção de variantes predefinidas de terminação de instruções, blocos, tipagem e arranjos. O conjunto de produções e as alternativas aceitas pelo analisador são definidos na implementação: o usuário não acrescenta novas produções nem redefine livremente a semântica das operações. Assim, o caráter dinâmico corresponde à aplicação de uma configuração a cada novo processamento do código, dentro dos limites suportados pelo núcleo.

No serviço Python, os testes utilizam SQLite em memória e não verificam migrações ou características específicas do PostgreSQL. A inspeção também identificou lacunas de autorização por objeto e caminhos de listagem sem paginação obrigatória, descritos na Seção~\ref{subsec:limites-seguranca-backend}. A aprovação da suíte e a filtragem de campos nas respostas não constituem garantia de segurança integral \cite{projeto2026backend,projeto2026backendtestes,silva2026verificacaobackend}.

Por fim, o interpretador executa o código intermediário por percurso direto das instruções, sem etapa de otimização. As técnicas discutidas no referencial teórico foram estudadas, mas não implementadas, uma vez que o porte dos programas produzidos em contexto introdutório não as torna necessárias.

\section{Encaminhamentos}

A etapa seguinte do Trabalho de Conclusão de Curso concentrará esforços no aprofundamento do núcleo de tradução e na conclusão do protocolo de avaliação. No plano técnico, prevê-se a ampliação da cobertura de testes, sobretudo quanto à robustez do analisador diante das variações configuráveis da linguagem, e a verificação sistemática da correção do código intermediário gerado. No plano metodológico, prevê-se a consolidação dos critérios que permitirão aferir a contribuição da solução.

Para o back-end, os próximos passos incluem revisar as permissões entre usuários e instituições, acrescentar cenários negativos para os vínculos exigidos, limitar as listagens e conferir os controles no ambiente PostgreSQL. A política de credenciais e tokens também requer revisão periódica. Esses encaminhamentos complementam o fluxo já implementado de execução automática dos testes como condição para deploy \cite{owasp2026authorization,owasp2023recursos,owasp2026password,projeto2026workflow}.

A aplicação desse protocolo junto a estudantes, com a coleta de evidências empíricas em ambiente educacional, é registrada como trabalho futuro. Sua condução exigirá, além do protocolo consolidado, a articulação com turmas em andamento e a observância dos procedimentos éticos aplicáveis a pesquisas que envolvem seres humanos.

Como desdobramentos adicionais, identificam-se a ampliação do subconjunto da linguagem suportado pelo interpretador, a investigação de mecanismos que permitam configurar também produções gramaticais, além dos parâmetros e das variantes predefinidas, e o estudo da geração de código para um alvo executável fora do ambiente da plataforma.
